% =====================================================================
%  Chapter 11 — The \idxw{Keyhole} and How It Is Constructed
% =====================================================================

\chapter{The Keyhole and How It Is Constructed}
\label{ch:keyhole}

The keyhole\idx{Keyhole} is the most efficient idea in planetary defence, and it is also
the one most often misdescribed. It is not a place on the Earth. It is a
narrow window on a plane, and an object passes through it or it does not.
This chapter constructs the window from first principles, because the
construction is what determines how much a physical measurement is worth.

\section{The target plane}

When an asteroid's orbit is propagated to a close encounter with a planet, the
trajectory is described in a plane perpendicular to the planet's approach
direction: the target plane\idx{Target plane}, or B-plane. Two coordinates on that plane locate
the impact parameter: the component along the direction of the planet's
motion, called the along-track coordinate\idx{Along-track coordinate} \(\zeta\), and the component across
it, called the cross-track coordinate \(S\).

The impact parameter\idx{Impact parameter} itself is the distance of the object's centre from the
planet's centre at the encounter epoch,
\begin{equation}
b = \sqrt{S^{2} + \zeta^{2}},
\label{eq:bplane}
\end{equation}
and a collision occurs when \(b\) is smaller than the sum of the two radii.

\begin{deeper}{Why the target plane and not the sky}
Because the planet moves. The distance at which an object passes the Earth
depends on where it is in the sky relative to the Earth's motion at the
moment of passage, not on where it is relative to the stars. An observer on
the Earth's surface, looking at the same object on two different nights, sees
it in the same place in the stars, and the planet's own motion carries the
impact parameter\idx{Impact parameter} around by a very large amount. The target plane\idx{Target plane} removes
that motion and leaves the geometry of the encounter itself.
\end{deeper}

\section{From a collision to a keyhole}

The set of positions on the target plane\idx{Target plane} that produce a collision is a
region of finite width, because the collision is not a point. The width of
that region is set by the sum of the radii and by the geometry of the
approach: at a grazing encounter the effective collision cross-section is
enlarged, and the region is elongated along the direction of the relative
velocity.

The region is called a keyhole\idx{Keyhole}, and its dimensions for a given encounter are
computable from the published encounter geometry. For the 2036 encounter of
Apophis with the Earth, the published keyhole\idx{Keyhole} width is 616.16 metres.

That is a very narrow window, and the narrowness is the entire point. A
keyhole that were a thousand kilometres wide would be a corridor; a keyhole
six hundred metres wide is a target.

\section{The probability, and what it is a function of}

The impact probability is the fraction of the probability distribution of the
along-track coordinate\idx{Along-track coordinate} that falls inside the keyhole. If the along-track
distribution is Gaussian with standard deviation \(\sigma_\zeta\), and the
keyhole is a strip of half-width \(w/2\) centred on \(\zeta_c\), then
\begin{equation}
P = \int_{\zeta_c - w/2}^{\zeta_c + w/2}
\frac{1}{\sigma_\zeta\sqrt{2\pi}}
\exp\!\left(-\frac{\zeta^{2}}{2\sigma_\zeta^{2}}\right)\,\mathrm{d}\zeta.
\label{eq:keyhole-prob}
\end{equation}

Three things follow immediately, and they are worth separating because they
are confused constantly.

The probability is a function of the \emph{width} of the keyhole, which is
fixed by the geometry of the encounter. No amount of characterisation changes
it.

The probability is a function of the \emph{position} of the keyhole relative
to the distribution, which is set by the object's physical properties. This
is the term that a characterisation campaign moves.

The probability is a function of \(\sigma_\zeta\), the width of the
along-track distribution, which is set by the astrometry.

\section{Reconstructing the published number}

The published 2036 impact probability for Apophis is \(7.07\times 10^{-9}\),
the keyhole width is 616.16 metres, and the keyhole centre sits 4.77 sigma
out in the along-track distribution. Those three numbers are not independent,
and the book's own generator reconstructs them from one another.

Take the along-track distribution to be Gaussian with standard deviation
\(\sigma_\zeta\). At a centre position of \(z = -4.77\) standard deviations,
the probability density per unit length is
\begin{equation}
\phi(-4.77) = \frac{1}{\sqrt{2\pi}}\exp\!\left(-\frac{4.77^{2}}{2}\right)
= 4.57\times 10^{-6},
\end{equation}
in units of inverse sigma. The keyhole's full width in units of sigma is
\(616.16/\sigma_\zeta\). The probability is the product, and setting it equal
to the published \(7.07\times 10^{-9}\) and solving gives
\begin{equation}
\sigma_\zeta = 398.5\ \mathrm{km}.
\label{eq:sigmazeta}
\end{equation}

The arithmetic closes: with that value of \(\sigma_\zeta\), the computed
probability is \(7.07\times 10^{-9}\), which is the published number. The
along-track uncertainty on the 2036 encounter is therefore about four
hundred kilometres, and the keyhole is six hundred metres wide --- roughly
one and a half parts in a thousand of the distribution's width, sitting
almost five sigma out in its tail.

\begin{table}[!htbp]
\centering
\small
\caption{What one standard deviation of the along-track distribution is
worth, at the 2036 keyhole.}
\label{tab:tail-sensitivity}
\begin{tabularx}{\textwidth}{@{}P{1.55in}P{1.01in}Y@{}}
\toprule
\rowcolor{panel}
\textbf{Shift of the keyhole centre} & \textbf{Probability factor} &
\textbf{Interpretation} \\
\midrule
0.1 per cent of \(\sigma\) & 1.00 & negligible \\
\rowcolor{panel}
1 per cent of \(\sigma\) & 1.05 & a few kilometres \\
10 per cent of \(\sigma\) & 1.61 & forty kilometres \\
\rowcolor{panel}
100 per cent of \(\sigma\) & 118 & four hundred kilometres \\
\bottomrule
\end{tabularx}
\end{table}

\tabref{tab:tail-sensitivity} is the magnification that the tail provides,
and it is derived from the Gaussian directly: the probability ratio between
two positions \(z_1\) and \(z_2\) in the tail is
\(\exp\left(z_{1}(z_{2}-z_{1})\right)\), and at \(z = 4.77\) the exponent is
large enough that a shift of one standard deviation changes the answer by a
factor of 118.

\section{What this chapter has established}

The keyhole is a narrow window in the tail of the along-track distribution.
Its width is geometric and fixed. Its position is physical, and the physical
assumptions that set its position are the same ones that set the object's
size and density. And because the window sits in the tail, a small change in
the physics is worth a large change in the probability.

The next three chapters work through what that physical position depends
on, and when it matters.

\begin{keytakeaways}
\item The impact probability is the fraction of the along-track distribution
that falls inside a window whose width is fixed by the encounter geometry.
\item The window's position relative to the distribution is physical, and it
is the term a characterisation campaign moves.
\item The published 2036 numbers are mutually consistent: a 616.16-metre
keyhole at 4.77 sigma with an along-track sigma of 398.5 kilometres
reproduces the published probability of \(7.07\times 10^{-9}\) exactly.
\item At that position in the tail, a shift of one standard deviation in the
keyhole centre changes the probability by a factor of 118.
\end{keytakeaways}
